\rok{Децембар 2025.}{C}

Први практични тест из Алгоритама и структура података

\zadatak{1}{Insertion Sort}
Написати Insertion Sort алгоритам за сортирање целобројног низа дужине $n$.

\textbf{Додатне напомене за решавање задатка:}
\begin{itemize}
	\item Улаз је несортиран низ целих бројева.
	\item Излаз је сортирани низ целих бројева.
	\item У template-у се налази класа \texttt{RandomArrayGenerator} коју треба искористити за генерисање низа случајних бројева.
	\item Креирати класу \texttt{InsertionSort} и у оквиру ње генерисати методу:
	\begin{Verbatim}[fontsize=\small]
public static void sort(long[] arr)
	\end{Verbatim}
\end{itemize}

\zadatak{2}{Кумулативна сума вредности левих суседа}
Написати алгоритам који модификује листу тако што вредност чвора замењује вредношћу суме свих елемената који се у листи налазе са његове леве стране.

\textbf{Додатне напомене за решавање задатка:}
\begin{itemize}
	\item Захтеван потпис методе: \texttt{public static void calculateCumulativeSum()}.
	\item Методу писати у оквиру класе \texttt{LinkedList}.
\end{itemize}

\textbf{Потпис методе:}
\begin{Verbatim}[fontsize=\small]
public static void calculateCumulativeSum()
\end{Verbatim}

\textbf{Пример:}
\begin{itemize}
	\item \textbf{Улаз:} \texttt{LinkedList: 1 -> 0 -> 6 -> 8 -> 4} \\
		  \textbf{Излаз:} \texttt{0 -> 1 -> 1 -> 7 -> 15}
	\item \textbf{Улаз:} \texttt{LinkedList: 1 -> 2 -> 4 -> 0 -> 6} \\
		  \textbf{Излаз:} \texttt{0 -> 1 -> 3 -> 7 -> 7}
\end{itemize}

\zadatak{3}{Towers of Hanoi}
Ханојске куле су класична математичка игра која се састоји од три штапа и одређеног броја дискова различитих величина који се могу низати на штапове. Игра почиње са свим дисковима наслаганим на једном штапу у растућем редоследу (најмањи на врху), формирајући облик конуса.

\textbf{Проблем:} Циљ је пребацити цео сет дискова са почетног штапа (Извор) на завршни штап (Циљ), користећи помоћни штап (Помоћни), поштујући следећа правила:
\begin{enumerate}
	\item Само један диск се може померати у датом тренутку.
	\item Сваки покрет се састоји од узимања горњег диска са једног од штапова и постављања на врх другог штапа.
	\item Златно правило: Ниједан диск не сме бити постављен изнад диска који је мањи од њега.
\end{enumerate}

\textbf{Додатни захтеви за решавање задатка:}
\begin{enumerate}
	\item Задатак решавати у оквиру методе:
	\begin{Verbatim}[fontsize=\small]
public void SolveHanoi(int n, Stack source, Stack destination, Stack auxiliary)
	\end{Verbatim}
	\item Симулација помоћу стекова:
	\begin{itemize}
		\item Сваки од три штапа (Извор, Циљ, Помоћни) мора бити представљен објектом класе \texttt{Stack}.
		\item Дискови су представљени целим бројевима (нпр. диск величине 3 је већи од диска величине 1).
		\item Операције померања дискова морају се вршити искључиво методама \texttt{push()} и \texttt{pop()}.
	\end{itemize}
	\item Резултат извршавања:
	\begin{itemize}
		\item Метода треба да исписује сваки корак (нпр. \texttt{"Померен диск 1 са штапа A на штап C"}).
		\item На крају извршавања, \texttt{source} и \texttt{auxiliary} стекови морају остати празни, а \texttt{destination} стек мора садржати све дискове у исправном редоследу.
	\end{itemize}
\end{enumerate}

\textbf{Пример:}

\textbf{Улаз:} \texttt{[6, 5, 4, 3, 2, 1]}

\begin{itemize}
	\item \textbf{Корак 0:} \\
		  \texttt{Source Stack \{6, 5, 4, 3, 2, 1\}} -- 1 је на врху стека \\
		  \texttt{Destination Stack \{\}} \\
		  \texttt{Auxiliary Stack \{\}}
	\item \textbf{Корак 1:} \\
		  \texttt{Source Stack \{6, 5, 4, 3, 2\}} -- 2 је на врху стека \\
		  \texttt{Destination Stack \{\}} \\
		  \texttt{Auxiliary Stack \{1\}} -- 1 је на врху стека
	\item \textbf{Корак 2:} \\
		  \texttt{Source Stack \{6, 5, 4, 3\}} -- 3 је на врху стека \\
		  \texttt{Destination Stack \{2\}} -- 2 је на врху стека \\
		  \texttt{Auxiliary Stack \{1\}} -- 1 је на врху стека
	\item \textbf{Корак 3:} \\
		  \texttt{Source Stack \{6, 5, 4, 3\}} -- 3 је на врху стека \\
		  \texttt{Destination Stack \{2, 1\}} -- 1 је на врху стека \\
		  \texttt{Auxiliary Stack \{\}}
	\item \textbf{Корак n:} \\
		  \texttt{Source Stack \{\}} \\
		  \texttt{Destination Stack \{6, 5, 4, 3, 2, 1\}} -- 1 је на врху стека \\
		  \texttt{Auxiliary Stack \{\}}
\end{itemize}

\sablon{Шаблон првог теста децембар 2025. група C}{2025/Decembar/GrupaC/AISP2025_Test1_GrupaC.linq}
